
\documentclass{article}
\usepackage{colortbl}
\usepackage{makecell}
\usepackage{multirow}
\usepackage{supertabular}

\begin{document}

\newcounter{utterance}

\centering \large Interaction Transcript for game `dond', experiment `semi', episode 16 with Qwen3.6{-}35B{-}A3B{-}FP8{-}without{-}reasoning.
\vspace{24pt}

{ \footnotesize  \setcounter{utterance}{1}
\setlength{\tabcolsep}{0pt}
\begin{supertabular}{c@{$\;$}|p{.15\linewidth}@{}p{.15\linewidth}p{.15\linewidth}p{.15\linewidth}p{.15\linewidth}p{.15\linewidth}}
    \# & $\;$A & \multicolumn{4}{c}{Game Master} & $\;\:$B\\
    \hline

    \theutterance \stepcounter{utterance}  
    & & \multicolumn{4}{p{0.6\linewidth}}{
        \cellcolor[rgb]{0.9,0.9,0.9}{
            \makecell[{{p{\linewidth}}}]{
                \texttt{\tiny{[P1$\langle$GM]}}
                \texttt{Pelaat yhteistoiminnallista neuvottelupeliä, jossa sinun on sovittava toisen pelaajan kanssa siitä, kuinka jakaa joukko kohteita teidän välillä.} \\
\\ 
\texttt{Säännöt:} \\
\texttt{(a) Sinä ja toinen pelaaja saatte joukon kohteita. Kummallekin annetaan myös salainen arvofunktio, joka kuvaa, kuinka paljon arvostit kunkin kohteen tyyppiä.} \\
\texttt{(b) Vaihdat viestejä toisen pelaajan kanssa sopimaan siitä, kuka saa mitkä kohteet. Jokainen voi lähettää enintään 5 viestiä tai lopettaa pelin aikaisin tekemällä salaisen ehdotuksen milloin tahansa.} \\
\texttt{(c) Sinua pyydetään lähettämään salainen ehdotus, joka osoittaa kohteet, jotka haluat, muotoiltuna hakasulkeiden sisällä seuraavasti: "{[}Ehdotus: <numero> <kohteen nimi>, <numero> <kohteen nimi>, <...>{]}"} \\
\texttt{(d) Jos ehdotuksesi ovat täydentäviä, eli molempien ehdotusten täyttämiseksi on tarpeeksi kohteita, jokainen pelaaja saa pistemäärän saamiensa kohteiden arvojen summan perusteella. Muussa tapauksessa molemmat saavat nolla pistettä.} \\
\texttt{(e) Tavoitteesi on maksimoida pistemääräsi ja toisen pelaajan pistemäärän summa.} \\
\\ 
\texttt{Aloitetaan.} \\
\\ 
\texttt{Saatavilla olevien kohteiden joukko on:} \\
\\ 
\texttt{2 banaanit, 2 potit, 1 kivi} \\
\\ 
\texttt{Salainen arvofunktiosi on:} \\
\\ 
\texttt{banaani: 4, potti: 1, kivi: 0} \\
\\ 
\texttt{TÄRKEÄÄ: Viestit, paitsi salaiset ehdotukset, lähetetään suoraan toiselle pelaajalle, joten älä sisällytä mitään vastausta sääntöihin tai tekstiä, joka ilmoittaa viestistäsi. Salaisesti ehdotuksen tekemiseksi käytä ilmoitettua muotoa. Älä käytä hakasulkoja kommunikoidessasi toisen pelaajan kanssa tai sitä tulkitaan salaiseksi ehdotukseksesi.} \\
            }
        }
    }
    & \\ \\

    \theutterance \stepcounter{utterance}  
    & \multicolumn{4}{p{0.6\linewidth}}{
        \cellcolor[rgb]{0.9,0.9,0.9}{
            \makecell[{{p{\linewidth}}}]{
                \texttt{\tiny{[P1$\rangle$GM]}}
                \texttt{Hei! Tässä on ehdotus jakamiseksi:} \\
\\ 
\texttt{{[}Salainen ehdotus: 2 banaania, 1 kivi{]}} \\
            }
        }
    }
    & & \\ \\

    \theutterance \stepcounter{utterance}  
    & & & \multicolumn{2}{p{0.3\linewidth}}{
        \cellcolor[rgb]{0.9,0.9,0.9}{
            \makecell[{{p{\linewidth}}}]{
                \texttt{\tiny{[GM$|$GM]}}
                \texttt{proposal does not start with "ehdotus:"} \\
            }
        }
    }
    & & \\ \\

\end{supertabular}
}

\end{document}
